\subsection{The actual dual half-level marking and its root representatives}
\label{subsec:ftb-actual-dual-half-marking}

We reconstruct the source's half-level linear algebra, since the
marking need not be an isomorphism modulo the whole integer level.
Let \(R\) be a commutative ring with an involution \(a\mapsto\bar a\),
let \(I,J\) be ideals and \(m>0\) an integer, and choose \(c\in R\)
such that
\[
 IJ\subset(m),\quad m\in I\cap J,\quad
 \bar I\subset J,\quad\bar J\subset I,\quad
                         c\in I,\quad c-1\in J.           \tag{DM1}
\]
Let \(A,P\) be \(R\)-modules, with \(P\) projective, and let
\(C\) and \(B\) be perfect integer-valued bilinear forms on
\(A\) and \(P\), respectively.
Perfect means that the maps to their full
\(\mathbb Z\)-duals are bijective.
Assume the adjoint rules
\[
 B(av,w)=B(v,\bar a w),\qquad
 C(ax,y)=C(x,\bar a y).                                   \tag{DM2}
\]
An \(R\)-linear marking \(f:A\to P\) has the source's two
half-level properties if
\[
 \begin{split}
 a\in IA,\ f(a)\in mP&\ \Longrightarrow\ a\in mA,\\
 v\in IP&\ \Longrightarrow\
        v=f(a)+mu\quad\text{for some }a\in IA,\ u\in P .
 \end{split}                                             \tag{DM3}
\]
These are the only marking injectivity and surjectivity
assumptions in the argument.

\begin{lemma}[A partial inverse and its adjoint]
\label{lem:ftb-dual-half-inverse}
Under (DM1)--(DM3), there is an \(R\)-linear map \(h:P\to A\)
with
\[
                    h(f(a))-a\in mA\qquad(a\in IA).
                                                               \tag{DM4}
\]
For each integer \(k\) there is an \(R\)-linear map
\(\delta:A\to P\) satisfying
\[
          m\mid kC(x,a)+B(\delta x,f(a))
                         \qquad(x\in A,\ a\in IA).        \tag{DM5}
\]
\end{lemma}
\begin{proof}
First, \(f(a)\in JP\) implies \(a\in JA\).
Indeed \(cf(a)\in mP\) by \(IJ\subset(m)\), and \(ca\in IA\).
The first clause of (DM3) gives \(ca=mb\).
Then \(a=mb-(c-1)a\in JA\).
Thus \(A/JA\to P/JP\) induced by \(f\) is injective.

The composite \(A\to P/JP\) is onto.
For \(v\in P\), apply the second clause of (DM3) to \(cv\)
to get \(cv=f(a)+mu\).
Then \(f(a)-v=(c-1)v-mu\in JP\).
Projectivity of \(P\) lifts its quotient map \(P\to P/JP\)
through this surjection \(A\to P/JP\), giving \(h:P\to A\)
with \(f(h(v))-v\in JP\).
For \(a\in IA\) the preceding injectivity test gives
\(e=h(f(a))-a\in JA\).
It also lies in \(IA\), since \(h,f\) are \(R\)-linear.
Consequently \(ce\in mA\) and \((c-1)e\in mA\);
subtracting gives \(e\in mA\), proving (DM4).
This is the source's projective-lifting construction, not
a choice of a whole-level inverse.

Perfectness of \(B\) supplies the unique adjoint \(h^\dagger\)
defined by
\[
                  B(h^\dagger x,v)=C(x,h(v)).
\]
It is additive by perfectness.
For \(a\in R\), (DM2) gives, for every \(v\),
\[
 C(ax,h(v))=C(x,\bar a h(v))
           =C(x,h(\bar a v))
           =B(a h^\dagger x,v),
\]
so it is \(R\)-linear.
Set \(\delta=-k h^\dagger\).
For \(a\in IA\), write \(h(f(a))=a+mb\) by (DM4).
Then
\[
 kC(x,a)+B(\delta x,f(a))=-kmC(x,b),
\]
which proves (DM5).
\end{proof}

\begin{lemma}[Half-level annihilators, dual surjectivity, and matching]
\label{lem:ftb-dual-half-kernel-match}
Keep (DM1)--(DM5), assume \(\gcd(k,m)=1\), and suppose \(B\)
is alternating.
Then
\[
 P=\delta(A)+IP,\qquad
             \delta(x)\in IP\ \Longleftrightarrow\ x\in IA.       \tag{DM6}
\]
If a second pair \(f',\delta'\) satisfies (DM3), (DM5) for
the same \(k\), and
\(\delta(x)-\delta'(x)\in IP\) for every \(x\), then
\[
                   f(a)-f'(a)\in mP\qquad(a\in IA).        \tag{DM7}
\]
\end{lemma}
\begin{proof}
We record the annihilator test actually used in the source.
For either of the perfect forms \(D\) on a module \(Q\),
\[
 \left[m\mid D(v,y)\text{ for all }y\in IQ\right]
                                  \Longrightarrow v\in IQ.      \tag{DM8}
\]
The hypothesis applied to \(cy\) says
\(m\mid D(\bar c v,y)\) for all \(y\in Q\).
The functional \(y\mapsto D(\bar c v,y)/m\) is integral and
\(\mathbb Z\)-linear.  Perfectness gives \(z\) with
\(\bar c v=mz\).
Thus \(v=mz-\overline{c-1}v\in IQ\).
We also use
\[
                         B(IP,IP)\subset m\mathbb Z,      \tag{DM9}
\]
and the corresponding statement for \(C\).
For a generator \(a u\) in the first \(I\)-module, move
\(a\) to the second variable using (DM2);
\(\bar a\,IP\subset JIP\subset mP\).
Extend by addition.  This proves (DM9) without an
integral basis diagonalizing \(R\).

If \(\delta x\in IP\), both \(\delta x\) and \(f(a)\) are
in \(IP\) for \(a\in IA\).
By (DM5), (DM9) and coprimality,
\(m\mid C(x,a)\) for all \(a\in IA\).
The test (DM8) gives \(x\in IA\).
The reverse implication follows from \(R\)-linearity.

For the first assertion of (DM6), take the adjoint
\(g:P\to A\) of \(f\), so \(C(gv,a)=B(v,f(a))\).
Choose integers \(p,q\) with \(pm+qk=1\), and put
\(x=-qgv\).
To show \(v-\delta x\in IP\), use (DM8).
Every \(y\in IP\) has the form \(f(a)+mu\) with \(a\in IA\).
The \(mu\) term pairs to a multiple of \(m\).
For the other term, (DM5) gives
\[
 B(\delta x,f(a))\equiv-kC(x,a)
                   =qkB(v,f(a))\pmod m.
\]
Hence \(B(v-\delta x,f(a))\equiv
(1-qk)B(v,f(a))=pmB(v,f(a))\equiv0\pmod m\).
This proves the required surjectivity.

Finally fix \(a\in IA\) and put \(d=f(a)-f'(a)\in IP\).
Subtract the two instances of (DM5).
The term \(B((\delta-\delta')x,f'(a))\) is divisible by
\(m\), by (DM9).  Thus \(m\mid B(\delta x,d)\).
Given any \(y\in P\), write \(y=\delta x+v\) with \(v\in IP\)
by (DM6).  Again (DM9) implies \(m\mid B(y,d)\).
Alternation gives \(m\mid B(d,y)\) for every \(y\).
Dividing this integer functional by \(m\) and using
perfectness gives \(d=mu\), proving (DM7).
\end{proof}

\begin{proposition}[The actual dual marking of the self-dual lattice]
\label{prop:ftb-actual-dual-half-marking}
Let \(E,H,D_w,l,X,L,\mathrm{mk},\zeta\) be the source data:
\(E\) is a CM number field, \(H\) its host, \(D_w\) its datum,
\(l\) its actual level, \(X\) its Hermitian data, \(L\) a
self-dual \(O_E\)-lattice in \(E^r\), \(\mathrm{mk}\) an
actual \(\mathrm{IsMark}(l,L)\), and \(\zeta\) a primitive
\(M=l.M\)-th root of unity.
Write
\[
             I=\mathrm{Mfr}(l),\qquad J=l.M_b.
\]
There is an \(O_E\)-linear
\(\delta:D_w.\mathrm{La}\to E^r\) with image in \(L\) and
the exact source dual-pairing relation
\[
 \zeta^{\,D_w.EF(x,c_{\rm Sep}a)}
   =\exp\!\left\{-2\pi i\,
       EF_X\bigl(\delta x,M^{-1}\mathrm{mk}(a)\bigr)\right\}
                         \qquad(a\in I D_w.\mathrm{La}). \tag{DM10}
\]
Every such \(\delta\) satisfies
\[
 L\subset\operatorname{im}\delta+IL,\qquad
 \delta a\in IL\ \Longleftrightarrow\
                         a\in I D_w.\mathrm{La}.          \tag{DM11}
\]
If \(\mathrm{mk}',\delta'\) is another such marking/dual pair
for the same data and \(\delta-\delta'\) has image in \(IL\),
then
\(\mathrm{mk}(a)-\mathrm{mk}'(a)\in ML\) for every
\(a\in I D_w.\mathrm{La}\).
For \(u\) in the actual unitary group of \(X\),
\[
  u\mathrm{mk},u\delta
\]
are the corresponding marking and dual marking on \(uL\).
In addition
\[
                      \bar I=J,\qquad IJ=(M).             \tag{DM12}
\]
\end{proposition}
\begin{proof}
Self-duality of \(L\) supplies a perfect integral alternating
form \(B\) with
\[
                    B(v,w)=EF_X(v,w)\quad(v,w\in L).
\]
This is the trace-form duality calculation:
the coefficient \(X.\kappa_0/2\ne0\), the Hermitian matrix
is nonsingular, and \(L=L^\sharp\); hence the map from \(L\)
to its full integer dual is bijective.
Its adjoint rule is
\(B(av,w)=B(v,\bar a w)\), obtained by applying the trace
identity and the host conjugation.
The datum form \(C=D_w.EF\) is perfect and has the same
adjoint rule.

The lattice \(L\) is a finite torsion-free \(O_E\)-module.
Since \(O_E\) is Dedekind it is flat, and its finite
presentation makes it projective.
We do not require it to be free over \(O_E\).
The actual level supplies (DM1) with
\(m=M,c=c_{\rm Sep}\), including (DM12).
The properties of \(\mathrm{IsMark}\), after factoring its
image through \(L\), are exactly (DM3).

Write \(\zeta=\exp(2\pi i k/M)\) with \(\gcd(k,M)=1\).
For \(a\in I D_w.\mathrm{La}\),
\((c_{\rm Sep}-1)a\in M D_w.\mathrm{La}\).
Consequently (DM10) is equivalent, by the primitive-root
test, to
\[
 M\mid kD_w.EF(x,a)+B(\delta x,\mathrm{mk}(a)).          \tag{DM13}
\]
Indeed \(EF_X(v,M^{-1}w)=M^{-1}EF_X(v,w)\);
the replacement of \(c_{\rm Sep}a\) by \(a\) changes the
integer exponent by a multiple of \(M\).
Apply Lemma~\ref{lem:ftb-dual-half-inverse} for existence,
and Lemma~\ref{lem:ftb-dual-half-kernel-match} for all
the asserted properties of an arbitrary dual marking.
The equality \(\bar I=J\) follows in both directions from
the involution and the defining conjugate ideal.
The product in (DM12) is the actual level product.
Finally an invertible \(u\) transports lattice membership
and the two marking properties; if \(u\) is unitary,
\(EF_X(uv,uw)=EF_X(v,w)\), so it transports (DM10) as well.
\end{proof}

\begin{lemma}[Actual root representatives with a fixed denominator power]
\label{lem:ftb-actual-half-level-roots}
For the same \(L,\mathrm{mk},l\), every \(N>0\) and every
number of rows \(n\) admit the root representatives in the
source predicate \(\mathrm{IsRoots}(l,L,\mathrm{mk},N)\).
More explicitly there are \(t\in O_E,j\ge0\) with
\[
       t\equiv1\pmod J,\qquad tJ^j\subset(2N),
\]
and, for each row class \(q\in
D_w.\mathrm{La}/J D_w.\mathrm{La}\), one can use
\[
            Y(q)=\frac{t}{M}
                       \mathrm{mk}(c_{\rm Sep}\operatorname{rep}(q)).
                                                               \tag{DM14}
\]
This is congruent modulo \(L\) to the indicated level lift,
and
\[
                         J^{j+1}Y(q)\subset2NL
\]
uniformly in \(q\).
The same \(t,j\) work for every row; matrices are assembled
row by row.
\end{lemma}
\begin{proof}
The ring \(O_E/(2N)\) is finite and Artinian.
The descending sequence
\(J^j(O_E/(2N))\) stabilizes.
For a stabilized finite module \(Q=JQ\), the determinant
trick applied to a finite set of generators gives
\(t\equiv1\pmod J\) with \(tQ=0\).
Thus \(tJ^j\subset(2N)\).
This is the source's Artinian-stabilization and
Nakayama/determinant-trick step.

For \(a=c_{\rm Sep}\operatorname{rep}(q)\) we have
\(a\in I D_w.\mathrm{La}\).
For every \(b\in J\), \(b\,\mathrm{mk}(a)\in ML\), by
\(JI\subset(M)\), and hence
\(bM^{-1}\mathrm{mk}(a)\in L\).
Since \(t-1\in J\), (DM14) has the required congruence
modulo \(L\).
The class represented by \(a\) is still \(q\), because
\(c_{\rm Sep}-1\in J\).
For \(b\in J^j,d\in J\), first write
\(dM^{-1}\mathrm{mk}(a)=z\in L\), and then
\(tb=2Nw\) with \(w\in O_E\).
Their product is \(2Nwz\in2NL\).
Additivity over sums of products proves
\(J^{j+1}Y(q)\subset2NL\), with the same exponent for all
classes and rows.
\end{proof}

\paragraph{Frozen-source correspondence.}
This is the actual producer
\texttt{src/Development/AbelianVarietyTheory/Proofs/%
SiegelThetaTransformation/DualHalfLevelMarkingExistsAndMatch.lean}
(\texttt{ca7df4a980bb13ba448a4268b4ac12dd635d7c9f}),
at the fixed source tag.
The displayed mechanism follows, in order,
\srcid{DualMark.exists_nsmul},
\srcid{smul_eq_nsmul}, \srcid{exists_adj},
\srcid{dvd_of_mem}, \srcid{mem_of_dvd},
\srcid{exists_inv}, \srcid{exists_dual},
\srcid{mem_of_dual_mem}, \srcid{exists_dual_sub_mem},
and \srcid{mark_eq_of_dual}.
The actual lattice specialization is
\srcid{exists_form}, \srcid{isDual_iff},
\srcid{conj_half}, \srcid{isMark_map},
\srcid{isDual_map}, and \srcid{dual_marking}.
The denominator-power construction is
\srcid{exists_pow_le} and \srcid{exists_roots}.
Its exported interface is
\texttt{DualHalfLevelMarkingExistsAndMatchStatement.lean}
(\texttt{7f5fac9f3799aaede0436a4225175cb22ce164ab}),
\srcid{GlobalTheory.T.dual_half_level_marking_exists_and_match}.
This supplies the actual higher-level dual marking retained
in the cover-Fourier construction, not only the downstream
CRT consequences of an assumed dual marking.
